\usepackage{etoolbox}
% A navigable outline in the reader's sidebar: present, expanded to the chapter
% level, and shown when the file opens. Deferred to the start of the document:
% pandoc loads hyperref AFTER this file, so calling \hypersetup here directly
% left it undefined and LaTeX typeset the options onto the first page.
\AtBeginDocument{\hypersetup{bookmarks=true, bookmarksopen=true,
  bookmarksopenlevel=1, bookmarksnumbered=false, pdfpagemode=UseOutlines}}
\usepackage{graphicx}
\usepackage{float}
\usepackage{needspace}
% No paragraph leaves its last line alone at the top of a page, or its first
% line alone at the foot of one.
\widowpenalty=10000
\clubpenalty=10000
\usepackage{newunicodechar}
\usepackage{titlesec}
\usepackage{tocloft}
\usepackage{array}
\usepackage{ragged2e}
\usepackage[most]{tcolorbox}
\usepackage{fvextra}
\usepackage{caption}

% ── Typefaces: the same pair the website uses ────────────────────────────────
% Body ACaslon Pro, headings Test National.
\setmainfont[Path=FONTSDIR/,
  UprightFont=ACaslonPro-Regular.otf,
  ItalicFont=ACaslonPro-Italic.otf,
  BoldFont=ACaslonPro-Semibold.otf,
  BoldItalicFont=ACaslonPro-SemiboldItalic.otf]{ACaslonPro}
\setsansfont[Path=FONTSDIR/,
  UprightFont=TestNational-Regular.otf,
  ItalicFont=TestNational-BookItalic.otf,
  BoldFont=TestNational-Semibold.otf,
  BoldItalicFont=TestNational-Semibold.otf]{TestNational}

% ── Glyphs the heading face does not have ────────────────────────────────────
% The Test National trial carries 68 glyphs: letters, digits and a little
% punctuation. It has no colon, slash, question mark, en dash or quotation
% marks, so headings silently lost those characters. The browser hides this by
% falling back per character to Helvetica Neue, which is what --sans says to do;
% XeLaTeX has no such fallback, so do it explicitly — and only inside a heading,
% or every quotation mark in the body would turn into Helvetica too.
% The seven characters the heading face lacks are wrapped in \headfallback by
% a pass over the generated LaTeX (step 4 of build.sh). newunicodechar cannot
% help here: it refuses plain ASCII, and ":", "/" and "?" are the common cases.
\newfontfamily\headfallback{Helvetica Neue}
\newfontfamily\monofont{Menlo}

\newfontfamily\cjkfallback{Arial Unicode MS}
\newunicodechar{質}{{\cjkfallback 質}}
\newunicodechar{問}{{\cjkfallback 問}}
\newunicodechar{问}{{\cjkfallback 问}}
\newunicodechar{̂}{{\cjkfallback ̂}}
\newunicodechar{κ}{{\cjkfallback κ}}
\newunicodechar{α}{{\cjkfallback α}}
\newunicodechar{В}{{\cjkfallback В}}
\newunicodechar{о}{{\cjkfallback о}}
\newunicodechar{п}{{\cjkfallback п}}
\newunicodechar{р}{{\cjkfallback р}}
\newunicodechar{с}{{\cjkfallback с}}
\newunicodechar{‐}{{\cjkfallback ‐}}
\newunicodechar{→}{{\cjkfallback →}}
\newunicodechar{↓}{{\cjkfallback ↓}}
\newunicodechar{↑}{{\cjkfallback ↑}}
% 𝕏 (U+1D54F) appears in a tweet title; only a maths face carries it.
\newfontfamily\mathletterfallback{STIX Two Math}
\newunicodechar{𝕏}{{\mathletterfallback 𝕏}}

% ── Headings: larger, sans, and flush to the top of the page ─────────────────
% titlesec's third argument is the space BEFORE the title; zero puts a chapter
% on the first line of its page. The fourth is the space after, which is what
% keeps the title clear of the text.
% [block], not [display]: the display shape reserves a line for the (empty)
% chapter label, which left 2cm of white above every chapter title.
% A chapter opens with its number on its own line, a size down, a clear gap,
% then the full title. The label is set per chapter by \chapterwithlabel,
% which headglyphs.py substitutes for \chapter; front matter, Parts and the
% Conclusion have no label and print the title alone.
\newcommand{\chapterlabel}{}
\newcommand{\chapterwithlabel}[3]{%
  \renewcommand{\chapterlabel}{#1}\chapter[#2]{#3}\renewcommand{\chapterlabel}{}}
\titleformat{\chapter}[display]
  {\normalfont\sffamily\bfseries\fontsize{27}{29}\selectfont\linespread{0.92}\selectfont\raggedright
   % pandoc sets secnumdepth to -\maxdimen, and titlesec drops the label slot
   % of an unnumbered heading entirely, so the "Chapter N" line is typeset here
   % in the format code instead. \titlespacing pulls every chapter up by 21pt
   % so an unlabelled title sits on the top margin; a labelled one gives that
   % back first, or the number line lands above the margin and is clipped.
   % \endgraf rather than \par: titleformat's arguments are not \long.
   \ifx\chapterlabel\empty\else
     \vspace*{21pt}{\fontsize{17}{20}\selectfont\chapterlabel\endgraf}\vspace{30pt}\fi}
  {}{0pt}{}
% The negative value cancels the leading that 1.4 line spacing adds to the
% title's own first line; measured so the cap height lands on the top margin.
\titlespacing*{\chapter}{0pt}{-21pt}{52pt}
\titleformat{\section}
  {\normalfont\sffamily\bfseries\fontsize{19}{21}\selectfont\linespread{0.92}\selectfont\raggedright}{}{0pt}{}
\titlespacing*{\section}{0pt}{22pt}{10pt}
\titleformat{\subsection}
  {\normalfont\sffamily\bfseries\fontsize{15.5}{17.5}\selectfont\linespread{0.92}\selectfont\raggedright}{}{0pt}{}
\titlespacing*{\subsection}{0pt}{18pt}{8pt}
\titleformat{\subsubsection}
  {\normalfont\sffamily\bfseries\fontsize{13.5}{15.5}\selectfont\linespread{0.92}\selectfont\raggedright}{}{0pt}{}
\titlespacing*{\subsubsection}{0pt}{16pt}{7pt}

% Contents entries are set in the heading face too, so they need the flag.
% tocloft typesets the "Contents" title itself, bypassing \chapter*, so the
% titlesec format above never reached it and it came out in the serif. Give
% it the chapter title's face, size and position.
\renewcommand{\cfttoctitlefont}{\normalfont\sffamily\bfseries\fontsize{27}{29}\selectfont}
\setlength{\cftbeforetoctitleskip}{-21pt}
\setlength{\cftaftertoctitleskip}{52pt}
\renewcommand{\cftchapfont}{\sffamily\bfseries}
\renewcommand{\cftsecfont}{\sffamily}
\renewcommand{\cftsubsecfont}{\sffamily}

% ── Tables ───────────────────────────────────────────────────────────────────
% Ragged right with hyphenation: a narrow column of prose justifies badly and
% long words were the reason cells collided in the first place.
\AtBeginEnvironment{longtable}{%
  \RaggedRight\setlength{\tabcolsep}{5pt}%
  \hyphenpenalty=10 \exhyphenpenalty=10 \emergencystretch=2em}
\AtBeginEnvironment{tabular}{%
  \RaggedRight\setlength{\tabcolsep}{5pt}%
  \hyphenpenalty=10 \exhyphenpenalty=10 \emergencystretch=2em}

% ── Callout / prompt boxes ───────────────────────────────────────────────────
% A one-cell table is a callout in this manuscript. The website draws it as a
% tinted panel; match that rather than setting it as prose between two rules.
\newtcolorbox{calloutbox}{
  breakable, enhanced, boxrule=0.4pt, arc=2pt,
  colback=black!3, colframe=black!18,
  left=10pt, right=10pt, top=8pt, bottom=8pt,
  before skip=12pt, after skip=12pt,
}
\newenvironment{callout}{\begin{calloutbox}}{\end{calloutbox}}


% ── Code blocks: a tinted box in a mono face, as on the website ──────────────
% Long prompt lines ran off the page, so break them anywhere rather than let
% them overflow the measure.
\fvset{breaklines=true, breakanywhere=true, fontsize=\small, commandchars=\\\{\}}
% Every code block, whatever fence it was written with, is put in this box by
% the filter, so all the prompts look alike and match the website's panel.
\newtcolorbox{promptbox}{
  breakable, enhanced, boxrule=0.4pt,
  colback=black!3, colframe=black!18, arc=2pt,
  left=8pt, right=8pt, top=7pt, bottom=7pt,
  before skip=12pt, after skip=6pt,
}
\ifdefined\Shaded\renewenvironment{Shaded}{\begin{promptbox}}{\end{promptbox}}
\else\newenvironment{Shaded}{\begin{promptbox}}{\end{promptbox}}\fi

% ── Epigraph at the head of a chapter ───────────────────────────────────────
% The website sets these in a mono face a little smaller than the body, with
% generous air above and below; match it.
\newenvironment{epigraph}
  {\par\vspace{18pt}\begingroup\monofont\small\color{black!70}\raggedleft}
  {\endgroup\par\vspace{22pt}}

% ── The "figures are interactive" notice ────────────────────────────────────
\newenvironment{chapternotice}
  {\par\vspace{14pt}\begingroup\itshape\small\color{black!65}\raggedright}
  {\endgroup\par\vspace{20pt}}

% ── Captions ────────────────────────────────────────────────────────────────
% Always centred, tight under the figure or table they belong to, with clear
% air before the text that follows.
\newenvironment{figcaption}
  {\par\nopagebreak\vspace{4pt}\begingroup\centering\small}
  {\par\endgroup\vspace{16pt}}
% Floats: never let a figure occupy a page LaTeX would rather fill with text,
% and allow a tall one its own page instead of dragging text down after it.
\setcounter{topnumber}{2}\setcounter{bottomnumber}{2}\setcounter{totalnumber}{3}
\renewcommand{\topfraction}{0.9}\renewcommand{\bottomfraction}{0.7}
\renewcommand{\textfraction}{0.1}\renewcommand{\floatpagefraction}{0.75}

% A display image sits close to its caption, so no trailing space of its own.
% Inside a float the image must leave room for its caption, or the float
% overruns the page by the caption's height ("Float too large for page by
% 52pt"). pandoc bounds images to \textheight; bound them a little tighter.
\newenvironment{figureblock}
  {\par\medskip\centering\setlength{\textheight}{0.8\textheight}}
  {\par\noindent}

